% Source: 02 - Tue Sep 29 2026.pdf, page 9. Content remains in source order.
\sourcepage{02}{9}
\sourceheading{EXAMPLES}
\begin{itemize}
\item Naturals must be encoded to the base $b\geq2$, $b=O(1)$.
\item A set of naturals is encoded by concatenating the encodings of its elements.
\end{itemize}
From now on we will ignore the details of the encoding and assume that a concise encoding is in use!

\textbf{NOTATION:} Given a set of abstract instances $\mathcal I$:
$\forall i\in\mathcal I$, $\langle i\rangle\in\bits^*$ denotes a concise encoding of $i$.
e.g., $\langle G=(V,E)\rangle$.

Standardization is now complete!
\[
\Pi_C:\bits^*\longrightarrow\{0,1\}.
\]
ANY concrete decision problem maps binary strings (instance encodings) into bits (1/0 encoding yes/no).
